\begin{pastexam}[label={px:ntu-106-esoe-ds-02}]{臺大工科海洋所 106 學年度 第 2 題}
（證明）

\textbf{(10 \%)} Let $f(n)=\sum_{i=1}^{n}\sqrt{i}=1+\sqrt{2}+\sqrt{3}+\cdots+\sqrt{n}$. By showing
both the upper and lower bounds, prove that $f(n)=\Theta(n\sqrt{n})$.
\end{pastexam}

\begin{pastsoln}
要找正的常數 $c_1, c_2$ 和 $n_0$，使得 $n \ge n_0$ 時
$c_1 n\sqrt{n} \le f(n) \le c_2 n\sqrt{n}$（資結書第~1.5~節 $\Theta$ 的定義）。 做法是資結書
第~A.2~節求和的老方法：上界把每一項放大成最大的那一項，下界丟掉前一半、把後一半縮小成其中最小
的那一項。

\textbf{上界。} $1 \le i \le n$ 時 $\sqrt{i} \le \sqrt{n}$，所以
\[
f(n) = \sum_{i=1}^{n} \sqrt{i} \le \sum_{i=1}^{n} \sqrt{n} = n\sqrt{n}.
\]
取 $c_2 = 1$，$f(n) = O(n\sqrt{n})$。

\textbf{下界。} 每一項都是正的，只留 $i \ge \lceil n/2 \rceil$ 的項，和只會變小。 這些項有
$n - \lceil n/2 \rceil + 1 = \lfloor n/2 \rfloor + 1 \ge n/2$ 項，每一項
$\sqrt{i} \ge \sqrt{\lceil n/2 \rceil} \ge \sqrt{n/2}$，所以
\[
f(n) \ge \sum_{i=\lceil n/2 \rceil}^{n} \sqrt{i}
\ge \frac{n}{2} \cdot \sqrt{\frac{n}{2}} = \frac{1}{2\sqrt{2}}\, n\sqrt{n}.
\]
取 $c_1 = \frac{1}{2\sqrt{2}}$，$f(n) = \Omega(n\sqrt{n})$。

兩個不等式對所有 $n \ge 1$ 都成立，取 $n_0 = 1$，得到
$\frac{1}{2\sqrt{2}}\, n\sqrt{n} \le f(n) \le n\sqrt{n}$，所以 $f(n) = \Theta(n\sqrt{n})$。
\hfill$\blacksquare$

另一個做法是用積分夾（資結書引理~A.6）：$\sqrt{x}$ 遞增，
$\int_0^n \sqrt{x}\,dx \le f(n) \le \int_1^{n+1} \sqrt{x}\,dx$，也就是
$\frac{2}{3} n^{3/2} \le f(n) \le \frac{2}{3}\bigl((n+1)^{3/2} - 1\bigr)$，兩邊都是
$\Theta(n^{3/2})$，還順便看出 $f(n) / (n\sqrt{n}) \to 2/3$。 題目要求「上下界都寫出來」，
兩種寫法都要把常數交代清楚。
\end{pastsoln}
